% Volume III: Attention and Publicity
% Part VI. Privacy as Boundary and Capacity
% Chapter 30: The Mosaic Problem
% Spine: proof-spines/ch030-spine.md

\chapter{The Mosaic Problem}
\label{ch:ch030}

The \textbf{mosaic problem} begins from the fact that individually harmless disclosures need not remain harmless in aggregate. If disclosures \(D_1,\ldots,D_n\) interact, aggregate privacy risk is generally nonadditive:
\[
R(D_1\cup\cdots\cup D_n) \neq \sum_i R(D_i).
\]
More fully,
\[
R(D_{1:n}) = \sum_i R(D_i) + \sum_{i<j} I_{ij} + \sum_{i<j<k} I_{ijk} + \cdots,
\]
\ToyModel\ where the synergistic inference terms \(I_{ij}, I_{ijk}, \ldots\) register what only combinations reveal.

The chapter uses this expansion to connect privacy loss to combinatorial causal inference. Item-by-item sensitivity tests miss the level at which harm emerges, because no single datum need contain the invasion later produced by linkage, accumulation, and recombination. Privacy can be lost not only by revelation, but by composition. The same resistance to recombination becomes concealment when aggregation rules are drawn so broadly that patterns of official or corporate power cannot be assembled either.
